\documentclass[11pt,oneside]{book}

\usepackage[margin=1.15in]{geometry}
\usepackage{amsmath,amssymb,amsthm,mathtools}
\usepackage{bm}
\usepackage{enumitem}
\usepackage{microtype}
\usepackage{booktabs}
\usepackage{tikz}
\usetikzlibrary{arrows.meta,positioning,calc}
\usepackage[hidelinks]{hyperref}
\usepackage{titlesec}

% ---- theorem environments (identical to the keystone, for series coherence) ----
\theoremstyle{definition}
\newtheorem{definition}{Definition}[chapter]
\newtheorem{axiom}{Axiom}[chapter]
\newtheorem{construction}{Construction}[chapter]
\theoremstyle{plain}
\newtheorem{theorem}{Theorem}[chapter]
\newtheorem{lemma}[theorem]{Lemma}
\newtheorem{proposition}[theorem]{Proposition}
\newtheorem{corollary}[theorem]{Corollary}
\theoremstyle{remark}
\newtheorem*{remark}{Remark}

% ---- notation macros (shared with projection_thesis.tex / 00_foundations.tex) ----
\newcommand{\Obs}{\mathcal{O}}                 % raw observation space
\newcommand{\Adm}{\mathcal{O}^{*}}             % admitted observation space
\newcommand{\Act}{\mathcal{A}}                 % action / artifact space
\newcommand{\muop}{\mu}                        % manufacturing morphism
\newcommand{\adm}{\alpha}                      % admission map
\newcommand{\chainH}{\mathsf{H}}               % hash
\newcommand{\Rfsl}{\bot}                        % refusal (NOT \Ref: reserved by hyperref)
\newcommand{\CL}{\kappa}                       % bounded-verifier context capacity
\newcommand{\Fitness}{\varphi}
\DeclareMathOperator{\dom}{dom}
\DeclareMathOperator{\im}{im}
\DeclareMathOperator{\id}{id}
\DeclareMathOperator{\supp}{supp}
\DeclareMathOperator*{\argmax}{arg\,max}
% ---- paper-01-local macros ----
\newcommand{\code}[1]{\texttt{#1}}
\newcommand{\Deny}{\mathbb{D}}                 % denial-word space
\newcommand{\Adml}{\mathsf{ADMITTED}}          % bottom / clean word
\newcommand{\firedmap}{\pi_{f}}                % fired-mask projection
\newcommand{\Catset}{\mathcal{C}}              % refusal-category set
\newcommand{\Scnset}{\mathcal{S}}              % refusal-scenario set
\DeclareMathOperator{\catop}{cat}
\DeclareMathOperator{\laneop}{lane}
\DeclareMathOperator{\scnop}{scn}

\title{\bfseries The Admission Monoid:\\[0.4em]
\large An Algebra of Refusal}
\author{The Chatman Equation Thesis Program \\ (Praxis / Capability Physics)}
\date{Genesis --- Pillar I, Paper 01}

\begin{document}

\frontmatter
\maketitle

\chapter*{Abstract}
\addcontentsline{toc}{chapter}{Abstract}
The foundations paper of this series established that admission cannot be a semantic
decision and must be a computable retraction, with \emph{refusal} promoted from an exception
to a first-class value \cite{chatman2025}. This paper develops the algebra of that value. We
take as our object the \emph{denial word} --- the machine-checkable record of \emph{why} an
observation was not admitted --- and prove that the space of denial words is a bounded,
commutative, idempotent monoid (a join-semilattice) whose identity is the clean word
$\Adml$ and whose product is componentwise disjunction, realized in \code{praxis} as the
\code{DenialPolarity} bitfield under \code{compose}. We show that obligations are
\emph{lane maps} into this monoid, that total denial is their join, and that admission is
\emph{antitone} in the obligation set: adding obligations can only shrink the admitted set
(monotonicity). We then formalize the eight-bucket \code{RefusalCategory} taxonomy as a
total classification $\catop:\Scnset\to\Catset$ and prove a \emph{mechanized totality
theorem}: $\catop$ is realized by a wildcard-free \code{match}, so it is compiler-certified
total, its image is exactly the seven inhabited buckets, and the eighth (\code{Reserved})
carries an empty preimage by design --- an honest slack coordinate, not an overclaim. We
identify precisely where the type system's totality guarantee ends and a runtime test must
continue: the lane decoder \code{scenario\_for\_denial\_lane} is a \emph{partial} map because
its domain is an open $2^{64}$ newtype rather than a closed enum. We prove that admission
composes across a manufacturing pipeline as the unique monoid homomorphism from the free
monoid on stages into the denial monoid, and that this composition is order- and
multiplicity-independent, so a pipeline admits exactly when every stage admits. Finally, we
refine admission at the logical layer: the \code{prolog8} kernel decides bounded Horn queries
by SLD-resolution and returns a \emph{proof-carrying trichotomy} --- an \code{Answered} query
ships a positive proof DAG, a \code{Denied} query ships a negative proof tree, and an
\code{Invalid} query ships a \code{RejectionCode}; we prove this trichotomy and show it embeds
back into the denial monoid. Every object is anchored to a line of the \code{praxis}
repository, so a reader can trace each theorem to the code that enforces it.

\tableofcontents

\mainmatter

\chapter{Introduction: The Algebra a Refusal Needs}
\label{ch:intro}

\section{From ``refusal is data'' to an algebra of refusal}

The prior published work \cite{chatman2025} established the governing identity
$\Act=\muop(\Adm)$ and demonstrated it at enterprise scale. The foundations paper of this
series (Paper~00) grounded it: it proved a specialization of Rice's theorem
\cite{rice1953} to observations, concluded that admission \emph{cannot} decide meaning and
\emph{must} retract onto a decidable sub-language, and promoted the \emph{refusal} $\Rfsl$
from a thrown exception to a returned value. The keystone thesis introduced, in a single
Construction, the \emph{denial monoid} $D=(\{0,1\}^{n},\lor,\bm 0)$ and one proposition
about it. This paper pays that Construction its full debt.

The reason a refusal needs an algebra, rather than merely a name, is scale. Among a fleet
of $\sim\!10^{12}$ agents, each with bounded context $\CL$ and none able to read another's
interior, a refusal must satisfy four demands simultaneously:

\begin{enumerate}[label=(D\arabic*),leftmargin=*]
\item \textbf{Composition.} A manufacture passes through many admission stages; each may
refuse; the fleet needs \emph{one} verdict that faithfully aggregates all of them.
\item \textbf{Classification.} A refusal must land in a \emph{finite, shared, legible}
vocabulary, so that a recipient who never saw the sender's interior can nonetheless act on
the \emph{kind} of refusal.
\item \textbf{Totality.} No refusal may fall through a crack: every failure mode must map to
a class, and the mapping must be checkable by machine, not by hope.
\item \textbf{Reason-carrying.} A refusal must carry \emph{why}; at the logical layer this
sharpens to carrying a \emph{proof}.
\end{enumerate}

Composition (D1) wants a monoid. Classification (D2) wants a surjection onto a small set.
Totality (D3) wants a mechanized exhaustiveness argument. Reason-carrying (D4) wants, at its
limit, proof-carrying evaluation. This paper supplies each in turn --- Chapters~\ref{ch:sl}
and~\ref{ch:pipe} for composition, Chapter~\ref{ch:tax} for classification and totality, and
Chapter~\ref{ch:proof} for proof-carrying refusal --- and anchors each to running code.

\section{First principle}

\begin{axiom}[Refusal is data]\label{ax:refusal}
There is a distinguished refusal value $\Rfsl$ and a space of \emph{reasons} $\mathcal{R}_{\!f}$.
A refusal is not the absence of an output but the pair $(\Rfsl,r)$ with $r\in\mathcal{R}_{\!f}$
a machine-checkable reason. Admission is therefore \emph{total} as a map into
$\Adm\cup(\{\Rfsl\}\times\mathcal{R}_{\!f})$: a correctly functioning admission always
returns, either an admitted observation or a refusal-with-reason.
\end{axiom}

Axiom~\ref{ax:refusal} is the design decision made concrete in \code{praxis} by the
\code{Andon} type, whose \code{Halted} variant carries both the unmet obligations and their
refusal-taxonomy classification:
\[
\code{Andon::Halted\{\ unmet:\ Vec<Obligation>,\ refusals:\ Vec<RefusalScenario>,\ at:\ u64\ \}}.
\]
The remainder of this paper is the study of the structure that $\mathcal{R}_{\!f}$ carries.
We fix, once, the code objects under study: the \code{DenialPolarity} word and its
\code{compose} operation (\code{crates/praxis-core}, re-exported from \code{bcinr-powl-receipt});
the \code{Obligation} enum and the \code{Judge}/\code{Admit} traits (\code{law.rs}); the
\code{RefusalScenario}/\code{RefusalCategory} taxonomy and its total maps (\code{refusal.rs});
and the \code{prolog8} \code{Kernel::query} logic gate (\code{ops.rs}). Each theorem below names
the artifact it governs.

\chapter{The Denial Semilattice}
\label{ch:sl}

\section{The abstract denial word}

\begin{definition}[Denial word space]\label{def:denialabstract}
Fix $n\in\mathbb{N}$ \emph{lanes}. The \emph{abstract denial word space} is
$\Deny_n=\{0,1\}^{n}$ with componentwise disjunction $\lor$ and the all-zero word $\bm 0$.
A lane $i$ is \emph{set} in $d$ when $d_i=1$; the \emph{support} $\supp d=\{i:d_i=1\}$ names
the fired lanes. Write $d\preceq d'$ for the componentwise order $\forall i\, (d_i\le d_i')$.
\end{definition}

\begin{proposition}[$\Deny_n$ is a bounded idempotent commutative monoid]\label{prop:monoid}
$(\Deny_n,\lor,\bm 0)$ is a commutative monoid in which every element is idempotent
($d\lor d=d$). It is the join-semilattice of the Boolean lattice $(2^{[n]},\subseteq)$ under
the support isomorphism $d\mapsto\supp d$; $\bm 0$ is its least element (bottom), and the
all-ones word $\bm 1$ its greatest. Consequently $d\lor d'$ is the least upper bound of
$d,d'$, and $\preceq$ is exactly the algebraic order $d\preceq d'\iff d\lor d'=d'$.
\end{proposition}

\begin{proof}
Disjunction is associative, commutative, and idempotent per coordinate, with unit $0$ per
coordinate; a finite product of such structures inherits all four laws, giving a bounded
commutative idempotent monoid. The map $d\mapsto\supp d$ carries $\lor$ to $\cup$ and $\bm 0$
to $\varnothing$, and is a bijection $\Deny_n\to 2^{[n]}$, hence a monoid isomorphism onto
$(2^{[n]},\cup,\varnothing)$. In any join-semilattice the identity
$d\preceq d'\iff d\lor d'=d'$ holds by antisymmetry of $\subseteq$; least/greatest elements
are $\varnothing$ and $[n]$, i.e.\ $\bm 0$ and $\bm 1$
\cite[Ch.~I]{birkhoff1967}. \end{proof}

Proposition~\ref{prop:monoid} is the algebra behind ``a refusal category is a join-support''
in the keystone: $\supp d(o)\subseteq[n]$ is a subset, and subsets join by union. We now show
that the \code{praxis} implementation realizes exactly this structure.

\section{The concrete word: \texttt{DenialPolarity}}

\begin{definition}[The polarity word]\label{def:denialcode}
In \code{praxis} the denial word is \code{DenialPolarity}, a \code{\#[repr(transparent)]}
newtype over \code{u64} carved into eight byte-lanes. The clean word is
$\Adml=\code{DenialPolarity(0)}$. Seven named nonzero constants each occupy one distinct byte
lane:
\[
\begin{array}{llll}
\code{PRECONDITION\_FAILED} & (\text{lane }1) & \code{RESOURCE\_EXHAUSTED} & (\text{lane }4)\\
\code{SLA\_BREACH} & (\text{lane }2) & \code{OBJECT\_LIFECYCLE\_VIOLATION} & (\text{lane }5)\\
\code{AUTHORIZATION\_DENIED} & (\text{lane }3) & \code{CONFORMANCE\_GATE\_FAILED} & (\text{lane }6)\\
& & \code{WATCHDOG\_DRAINED} & (\text{lane }7).
\end{array}
\]
The product is \code{compose}$(a,b)=\code{DenialPolarity}(a.0\mathbin{|}b.0)$, bitwise
\textsc{or}; the admission predicate is \code{is\_admitted}$(d)\iff d.0=0$.
\end{definition}

\begin{proposition}[The polarity word is the denial semilattice]\label{prop:code_is_sl}
Let $\Deny=(\{\code{DenialPolarity}\},\code{compose},\Adml)$. Then $\Deny$ is a bounded
commutative idempotent monoid, and \code{is\_admitted} is exactly the predicate ``is the
bottom element.'' The submonoid $\langle L\rangle$ generated by the seven named nonzero
constants together with $\Adml$ is isomorphic to $\Deny_7$ of
Definition~\ref{def:denialabstract} via ``which named lanes are nonzero.''
\end{proposition}

\begin{proof}
Bitwise \textsc{or} on \code{u64} is associative, commutative, idempotent, with identity the
zero word; \code{compose} inherits these, so $\Deny$ is a bounded commutative idempotent
monoid whose bottom is $\Adml$, and \code{is\_admitted}$(d)\iff d.0=0\iff d=\Adml$. Each named
constant $c_i$ ($1\le i\le 7$) has $c_i.0=\code{0xFF}\ll 8i$ occupying lane $i$ alone, so for
any subset $S\subseteq\{1,\dots,7\}$ the composite $\bigvee_{i\in S}c_i$ has byte lane $i$
nonzero iff $i\in S$. The map $\bigvee_{i\in S}c_i\mapsto \bm 1_S\in\Deny_7$ (the indicator of
$S$) is therefore a well-defined bijection $\langle L\rangle\to\Deny_7$ carrying \code{compose}
to $\lor$ and $\Adml$ to $\bm 0$: a monoid isomorphism. \end{proof}

\begin{remark}
Definition~\ref{def:denialcode} uses eight byte-lanes but only seven carry a named denial;
lane $0$ is the region of the clean word and is never set by a denial constant. Thus the
\emph{word} has seven informative lanes, while the \emph{category} taxonomy of
Chapter~\ref{ch:tax} has eight buckets, the eighth (\code{Reserved}) being an uninhabited
slack coordinate. The two ``eights'' are the byte width, not a claim of eight active denials;
we are careful never to conflate them.
\end{remark}

\section{The fired-mask projection}

The keystone's agent byte and its planetary admission sweep require collapsing a wide,
possibly multi-lane word onto a fixed-width bit vector without losing \emph{which} lanes
fired. This is \code{to\_fired\_mask}.

\begin{definition}[Fired-mask]\label{def:fired}
$\firedmap:\Deny\to\{0,1\}^{8}$ sends a word $d$ to the bit vector whose bit $j$ equals the
nonzero-indicator of byte lane $j$ of $d$: $\firedmap(d)_j=\mathbf 1[(d.0\gg 8j)\ \&\ \code{0xFF}\ne 0]$.
The \code{praxis} realization computes this branchlessly via the identity
$\mathbf 1[x\ne 0]=(x\mathbin{|}(-x))\gg 63$ on each lane.
\end{definition}

\begin{theorem}[Fired-mask is a semilattice homomorphism]\label{thm:firedhom}
$\firedmap$ is a homomorphism of bounded commutative idempotent monoids:
$\firedmap(\Adml)=\bm 0$ and, for all $a,b\in\Deny$,
\[
\firedmap\bigl(\code{compose}(a,b)\bigr)=\firedmap(a)\lor\firedmap(b).
\]
Restricted to $\langle L\rangle$ it is a monoid isomorphism onto the seven-bit image
$\{0\}\times\{0,1\}^{7}\subseteq\{0,1\}^{8}$.
\end{theorem}

\begin{proof}
$\firedmap(\Adml)=\bm 0$ since every lane of the zero word is zero. Byte lane $j$ of
$\code{compose}(a,b)=a.0\mathbin{|}b.0$ is $a_j\mathbin{|}b_j$ (bitwise \textsc{or} acts
lane-locally), and for byte values $\mathbf 1[a_j\mathbin{|}b_j\ne 0]
=\mathbf 1[a_j\ne 0]\lor\mathbf 1[b_j\ne 0]$ because a disjunction of bytes is nonzero iff some
operand is; hence $\firedmap(a\mathbin{|}b)_j=\firedmap(a)_j\lor\firedmap(b)_j$ for every $j$,
which is the claimed homomorphism law. On $\langle L\rangle$, Proposition~\ref{prop:code_is_sl}
identifies elements with subsets $S\subseteq\{1,\dots,7\}$, and $\firedmap$ sends such an
element to the indicator of $S$ in bits $1\ldots7$ (bit $0$ always zero); this is a bijection
onto $\{0\}\times\{0,1\}^{7}$ preserving join and bottom. \end{proof}

Theorem~\ref{thm:firedhom} is the projection principle at the algebra layer: an arbitrarily
composed denial word --- the aggregate of an entire pipeline (Chapter~\ref{ch:pipe}) ---
projects losslessly, as far as \emph{which lanes fired}, onto at most eight bits. This is the
homomorphism that makes the keystone's word-parallel fleet-admission sweep sound
(Proposition~\ref{prop:fleet}).

\chapter{Obligations as Lane Maps}
\label{ch:ob}

\section{Obligations and total denial}

\begin{definition}[Obligation and lane map]\label{def:ob}
An \emph{obligation} is a first-class, hashable value the payload must satisfy before
admission. In \code{praxis} the \code{Obligation} enum has exactly three kinds:
\code{Precondition\{predicate\_id, params\_hash\}}, \code{BlockingConstraint\{reason\}}, and
\code{EvidenceRequired\{evidence\_type\}}. Each obligation $g$ induces a \emph{lane map}
$\delta_g:\Obs\to\Deny$ with
\[
\delta_g(o)=\begin{cases}\Adml & \text{if }g\text{ is satisfied by }o,\\[2pt]
\ell(g) & \text{otherwise,}\end{cases}
\]
where $\ell(g)\in\Deny$ is the lane $g$ fires on failure. For an obligation \emph{set}
$G=\{g_1,\dots,g_m\}$ the \emph{total denial} is the join
\[
d_G(o)\;=\;\code{compose\_denials}\bigl(\{\delta_{g_i}(o):g_i\text{ unmet}\}\bigr)
\;=\;\bigvee_{i=1}^{m}\delta_{g_i}(o).
\]
\end{definition}

The right-hand fold is literally \code{compose\_denials} in \code{refusal.rs}: it maps each
refusal scenario to its lane via \code{denial\_lane} and folds with \code{compose} from the
seed $\Adml$. Because $\Adml$ is the monoid identity (Proposition~\ref{prop:code_is_sl}), an
empty set of unmet obligations composes to $\Adml$ --- the admitted word --- which is exactly
the documented behavior (``empty input composes to \code{ADMITTED}'').

\begin{proposition}[Bottom characterizes admission]\label{prop:bottom}
$\adm_G(o)\ne\Rfsl\iff d_G(o)=\Adml\iff$ every $g_i\in G$ is satisfied by $o$.
\end{proposition}

\begin{proof}
$d_G(o)=\bigvee_i\delta_{g_i}(o)=\Adml$ iff every summand is $\Adml$ (a join of elements of a
join-semilattice equals the bottom iff each element is the bottom, by
Proposition~\ref{prop:monoid}), i.e.\ iff no obligation is unmet. Admission returns
non-refusal exactly in that case (Axiom~\ref{ax:refusal}). \end{proof}

\section{Monotonicity of admission}

\begin{theorem}[Admission is antitone in the obligation set]\label{thm:mono}
Let $G\subseteq G'$ be obligation sets. Then for every observation $o$,
\[
d_G(o)\;\preceq\;d_{G'}(o),\qquad\text{hence}\qquad
\Adm_{G'}\;\subseteq\;\Adm_{G},
\]
where $\Adm_G=\{o:d_G(o)=\Adml\}$ is the admitted set under $G$. Adding obligations can only
enlarge denial and shrink admission; it can never admit an observation a smaller obligation
set refused.
\end{theorem}

\begin{proof}
$d_{G'}(o)=\bigvee_{g\in G'}\delta_g(o)=\Bigl(\bigvee_{g\in G}\delta_g(o)\Bigr)\lor
\Bigl(\bigvee_{g\in G'\setminus G}\delta_g(o)\Bigr)=d_G(o)\lor d_{G'\setminus G}(o)\succeq d_G(o)$,
since in a join-semilattice $x\preceq x\lor y$ always (Proposition~\ref{prop:monoid}). If
$o\in\Adm_{G'}$ then $d_{G'}(o)=\Adml$, the bottom; from $d_G(o)\preceq d_{G'}(o)=\Adml$ and
minimality of $\Adml$ we get $d_G(o)=\Adml$, i.e.\ $o\in\Adm_G$. \end{proof}

\begin{corollary}[Admission is an antitone map of lattices]\label{cor:galois}
The assignment $G\mapsto\Adm_G$ is an order-reversing map from the lattice of obligation sets
$(2^{\mathcal{G}},\subseteq)$ to the lattice of admitted subsets $(2^{\Obs},\subseteq)$:
$G\subseteq G'\Rightarrow\Adm_{G'}\subseteq\Adm_G$. In particular the empty obligation set
admits everything ($\Adm_\varnothing=\Obs$), and enlarging obligations monotonically tightens
the gate.
\end{corollary}

\begin{proof}
Order-reversal is Theorem~\ref{thm:mono}. With $G=\varnothing$ the join is empty, so
$d_\varnothing(o)=\Adml$ for all $o$ and $\Adm_\varnothing=\Obs$. \end{proof}

\section{The obligation-to-scenario map is total}

\begin{proposition}[Obligation classification is compiler-certified total]\label{prop:obtotal}
The map $\scnop:\code{Obligation}\to\Scnset$ implemented by \code{From<\&Obligation> for
RefusalScenario} is a total function realized by a wildcard-free \code{match} over the three
\code{Obligation} kinds:
\[
\code{Precondition}\mapsto\code{UnsatisfiedPrecondition},\quad
\code{BlockingConstraint}\mapsto\code{BlockingConstraint},\quad
\code{EvidenceRequired}\mapsto\code{MissingEvidence}.
\]
Because the \code{match} has no wildcard arm, adding a fourth \code{Obligation} kind without
extending $\scnop$ is a compile error; totality is therefore a static property of the program,
not a runtime hope.
\end{proposition}

\begin{proof}
A Rust \code{match} on a closed enum type-checks only if every constructor is covered or a
wildcard is present; the cited \code{match} covers all three constructors and uses no
wildcard, so the compiler certifies exhaustiveness (this is the mechanized-totality mechanism
formalized in Paper~00). Each arm returns a specific \code{RefusalScenario}, so $\scnop$ is a
total function. \end{proof}

Proposition~\ref{prop:obtotal} is the first half of the totality argument: every obligation
\emph{failure} becomes a scenario. The second half --- every scenario becomes a category ---
is Chapter~\ref{ch:tax}.

\chapter{The Eight-Bucket Taxonomy as a Surjection}
\label{ch:tax}

\section{Categories and scenarios}

\begin{definition}[The taxonomy]\label{def:tax}
The \emph{category set} is the eight-element enum
\[
\Catset=\{\code{Identity},\code{Capacity},\code{Topology},\code{Temporal},
\code{Lifecycle},\code{Authorization},\code{Prerequisites},\code{Reserved}\}.
\]
The \emph{scenario set} $\Scnset$ is the thirteen-variant \code{RefusalScenario} enum,
partitioned by origin into: three \emph{obligation-driven} variants
(\code{BlockingConstraint}, \code{MissingEvidence}, \code{UnsatisfiedPrecondition}); seven
\emph{denial-lane} variants, one per named nonzero lane of Definition~\ref{def:denialcode}
(\code{WatchdogDrained}, \code{PreconditionFailed}, \code{SlaBreach}, \code{AuthorizationDenied},
\code{ResourceExhausted}, \code{ObjectLifecycleViolation}, \code{ConformanceGateFailed}); and
three \emph{logic/andon} variants (\code{KernelDenied}, \code{KernelInvalid},
\code{AndonInvariantViolated}). The \emph{category map} $\catop:\Scnset\to\Catset$ is
\code{RefusalScenario::category}.
\end{definition}

\begin{center}
\small
\begin{tabular}{lll}
\toprule
Scenario & $\catop(\cdot)$ & origin\\
\midrule
\code{KernelInvalid} & \code{Identity} & logic\\
\code{ResourceExhausted} & \code{Capacity} & denial lane\\
\code{ConformanceGateFailed} & \code{Topology} & denial lane\\
\code{AndonInvariantViolated} & \code{Topology} & andon\\
\code{WatchdogDrained}, \code{SlaBreach} & \code{Temporal} & denial lanes\\
\code{BlockingConstraint}, \code{ObjectLifecycleViolation} & \code{Lifecycle} & obligation / lane\\
\code{AuthorizationDenied}, \code{KernelDenied} & \code{Authorization} & lane / logic\\
\code{MissingEvidence}, \code{UnsatisfiedPrecondition}, \code{PreconditionFailed} & \code{Prerequisites} & obligation / lane\\
\midrule
\emph{(none)} & \code{Reserved} & ---\\
\bottomrule
\end{tabular}
\end{center}

\section{The totality theorem}

\begin{theorem}[Mechanized totality of classification]\label{thm:total}
The category map $\catop:\Scnset\to\Catset$ satisfies:
\begin{enumerate}[label=(\roman*),leftmargin=*]
\item \textbf{Totality.} $\catop$ is a total function: every one of the thirteen scenarios has
exactly one category.
\item \textbf{Mechanization.} $\catop$ is realized by a wildcard-free \code{match}, so its
totality is compiler-certified: adding a scenario variant without extending $\catop$ fails to
compile. (The \code{praxis} test suite additionally pins each variant's value.)
\item \textbf{Image.} $\im\catop=\Catset\setminus\{\code{Reserved}\}$: exactly the seven
buckets $\{\code{Identity},\code{Capacity},\code{Topology},\code{Temporal},\code{Lifecycle},
\code{Authorization},\code{Prerequisites}\}$ are inhabited, and \code{Reserved} has empty
preimage.
\end{enumerate}
Consequently $\catop$ is surjective onto the inhabited sub-taxonomy but is \emph{not}
surjective onto $\Catset$; the deficiency $\{\code{Reserved}\}$ is exactly one bucket.
\end{theorem}

\begin{proof}
(i)--(ii) A Rust \code{match} on the closed \code{RefusalScenario} enum type-checks only under
exhaustiveness; the \code{category} method matches all thirteen constructors with no wildcard
arm, so the compiler certifies that every scenario is assigned, and each arm returns a single
\code{RefusalCategory}, making $\catop$ total and single-valued. (iii) Reading the thirteen
arms (Definition~\ref{def:tax} table) exhibits a preimage for each of the seven listed buckets
--- e.g.\ $\code{KernelInvalid}\in\catop^{-1}(\code{Identity})$,
$\code{ResourceExhausted}\in\catop^{-1}(\code{Capacity})$, and so on down the table --- so all
seven are in the image; and no arm returns \code{Reserved}, so $\catop^{-1}(\code{Reserved})
=\varnothing$. Hence $\im\catop$ is precisely the seven-element complement of
$\{\code{Reserved}\}$. \end{proof}

\begin{remark}[Honesty of the reserved bucket]
Theorem~\ref{thm:total}(iii) is deliberately \emph{not} a claim of $8/8$ surjectivity. The
codomain names a bucket the mechanism does not yet inhabit, and the mathematics records that
faithfully rather than inflating coverage. This is the doctrine's antibody clause at the type
level: a reserved coordinate is declared, unused, and provably unused, so no reader can be
misled into thinking eight kinds of refusal are in force when seven are. When a future
enforcement lane populates \code{Reserved}, the theorem's image statement is what must be
re-verified.
\end{remark}

\section{Where the type guarantee ends: the partial lane decoder}

Not every total map in this taxonomy is a wildcard-free \code{match}, and the exception is
instructive: it marks the precise boundary between what the type system certifies and what a
runtime test must certify.

\begin{definition}[Single-lane set and the lane decoder]\label{def:decoder}
Let $L=\{\Adml,c_1,\dots,c_7\}\subseteq\Deny$ be the clean word together with the seven named
single-lane constants. The \emph{lane decoder}
$\scnop_\ell:\Deny\rightharpoonup\Scnset$ is \code{scenario\_for\_denial\_lane}: it returns
$\code{None}$ on $\Adml$, returns the matching denial-lane scenario on each $c_i$, and returns
$\code{None}$ on every word not in $L$.
\end{definition}

\begin{proposition}[The lane decoder cannot be a total \texttt{match}, and why]\label{prop:partial}
$\scnop_\ell$ is a partial map, definite (as \code{Some}) only on $L\setminus\{\Adml\}$. It is
\emph{not} realizable as a wildcard-free exhaustive \code{match}, because its domain $\Deny$ is
a \code{\#[repr(transparent)]} newtype over \code{u64} --- an open value space of cardinality
$2^{64}$, not a closed enum --- so the compiler cannot enumerate its cases. In particular any
composed multi-lane word $d=c_i\lor c_j\notin L$ (a legitimate element produced by
\code{compose}) has $\scnop_\ell(d)=\code{None}$: the decoder recognizes single lanes only.
Exhaustiveness over the seven single lanes is therefore a \emph{runtime test obligation}
(discharged by the \code{category\_mapping\_is\_total\_over\_denial\_lanes} test), not a static
type obligation.
\end{proposition}

\begin{proof}
Rust's exhaustiveness checker operates on closed sum types; a newtype over \code{u64} exposes
$2^{64}$ inhabitants with no finite constructor set, so no wildcard-free \code{match} on its
value is well-typed --- the implementation is necessarily an \code{if}/\code{else} chain over
the eight known constants with a \code{None} fallback. On $\Adml$ it returns \code{None} by the
first guard; on each $c_i$ it returns the corresponding \code{Some} scenario; on any other word
(including every $c_i\lor c_j$ with $i\ne j$, which lies outside $L$) it falls through to
\code{None}. Thus $\dom^{+}\scnop_\ell=L\setminus\{\Adml\}$ where $\dom^{+}$ denotes the
\code{Some}-domain. \end{proof}

\begin{proposition}[The lane encoder is total and a section]\label{prop:section}
The \emph{lane encoder} $\laneop:\Scnset\to\Deny$ implemented by \code{denial\_lane} is a total
function realized by a wildcard-free \code{match} over all thirteen scenarios. Restricted to
the seven denial-lane scenarios $\Scnset_\ell$, it and $\scnop_\ell$ form a section--retraction
pair:
\[
\scnop_\ell\circ\laneop=\id_{\Scnset_\ell}\quad\text{and}\quad
\laneop\circ\scnop_\ell=\id_{L\setminus\{\Adml\}},
\]
so the seven single lanes and the seven denial-lane scenarios are in bijection. The remaining
six scenarios (obligation-driven and logic/andon) are mapped by $\laneop$ onto \emph{closest-fit}
lanes and are therefore not in the range of that bijection.
\end{proposition}

\begin{proof}
$\laneop$ matches all thirteen constructors with no wildcard, hence is total (same mechanism as
Theorem~\ref{thm:total}). For each of the seven denial-lane scenarios $s$, $\laneop(s)=c_i$ is
its own lane, and $\scnop_\ell(c_i)=s$ by Definition~\ref{def:decoder}, giving
$\scnop_\ell(\laneop(s))=s$; conversely $\laneop(\scnop_\ell(c_i))=\laneop(s)=c_i$. Both
round-trips are the identity on their respective seven-element sets, the round-trip pinned by
the cited test. The six non-lane scenarios map by $\laneop$ onto lanes already claimed by the
seven (e.g.\ \code{MissingEvidence}$\mapsto\code{PRECONDITION\_FAILED}$), so $\laneop$ is not
injective overall and those six are outside the bijection. \end{proof}

\begin{corollary}[The assembled totality statement]\label{cor:assembled}
Composing the total maps of Propositions~\ref{prop:obtotal} and~\ref{prop:section} with the
total classification of Theorem~\ref{thm:total}: \emph{every unmet obligation} (via
$\catop\circ\scnop$) and \emph{every non-}$\Adml$\emph{ single denial lane} (via
$\catop\circ\scnop_\ell$ on $L\setminus\{\Adml\}$) maps to exactly one of the seven inhabited
categories. No obligation failure and no fired single lane is left unclassified.
\end{corollary}

\begin{proof}
Immediate by composition of the three total (respectively \code{Some}-total) maps, each
single-valued into its codomain, landing in $\im\catop=\Catset\setminus\{\code{Reserved}\}$
(Theorem~\ref{thm:total}). \end{proof}

Corollary~\ref{cor:assembled} is precisely the totality claim demanded of a fleet vocabulary
(D3): a refusal never falls through a crack, and its class is legible to any recipient who
holds the eight-element enum, regardless of the sender's interior.

\chapter{Composition Across a Pipeline as a Monoid Action}
\label{ch:pipe}

\section{Pipelines and their aggregate denial}

A manufacture rarely faces a single gate. It traverses a \emph{pipeline}: \code{Judge}
evaluates obligations; \code{Admit} may deny; a \code{prolog8} kernel gate may refuse a logical
side-condition (Chapter~\ref{ch:proof}); an optional andon second-gate ring may fire. Each
stage emits a denial word; the fleet needs one.

\begin{definition}[Pipeline and stage denial]\label{def:pipeline}
Let $\mathsf{Stage}$ be the set of admission stages. Fix an observation $o$ and let
$\varphi_o:\mathsf{Stage}\to\Deny$ send a stage to the denial word it emits on $o$
(via its scenarios and \code{denial\_lane}). A \emph{pipeline} is a finite sequence
$w=s_1 s_2\cdots s_k\in\mathsf{Stage}^{*}$, an element of the free monoid on $\mathsf{Stage}$
under concatenation with unit the empty sequence $\varepsilon$. Its \emph{aggregate denial} is
\[
\Phi_o(w)\;=\;\code{compose\_denials}\bigl(\varphi_o(s_1),\dots,\varphi_o(s_k)\bigr)
\;=\;\bigvee_{i=1}^{k}\varphi_o(s_i),\qquad \Phi_o(\varepsilon)=\Adml.
\]
\end{definition}

\begin{theorem}[Pipeline denial is the free-monoid homomorphism]\label{thm:freehom}
$\Phi_o:(\mathsf{Stage}^{*},\cdot,\varepsilon)\to(\Deny,\code{compose},\Adml)$ is the unique
monoid homomorphism extending $\varphi_o$ along the inclusion
$\mathsf{Stage}\hookrightarrow\mathsf{Stage}^{*}$. Because the target is commutative and
idempotent, $\Phi_o$ factors through the finite join-semilattice generated by
$\{\varphi_o(s):s\in\mathsf{Stage}\}$; hence $\Phi_o(w)$ depends only on the \emph{set} of
distinct stage-denials occurring in $w$ --- it is invariant under reordering the stages and
under repeating them.
\end{theorem}

\begin{proof}
$\Phi_o(\varepsilon)=\Adml$ and $\Phi_o(w\cdot w')=\bigvee_{s\in w}\varphi_o(s)\lor
\bigvee_{s\in w'}\varphi_o(s)=\Phi_o(w)\lor\Phi_o(w')=\code{compose}(\Phi_o(w),\Phi_o(w'))$
by associativity of $\lor$; so $\Phi_o$ is a monoid homomorphism, and it agrees with
$\varphi_o$ on length-one words. Uniqueness is the universal property of the free monoid: any
homomorphism agreeing with $\varphi_o$ on generators is determined on all of
$\mathsf{Stage}^{*}$ by the homomorphism laws. For the factorization, commutativity gives
order-independence ($a\lor b=b\lor a$) and idempotence gives multiplicity-independence
($a\lor a=a$), so $\Phi_o(w)=\bigvee_{s\in\mathrm{set}(w)}\varphi_o(s)$ depends only on the
underlying set of distinct denials. \end{proof}

\begin{corollary}[A pipeline admits iff every stage admits]\label{cor:allstages}
$\Phi_o(w)=\Adml\iff\varphi_o(s_i)=\Adml$ for every stage $s_i$ in $w$. A single refusing stage
refuses the whole pipeline, and the aggregate word records \emph{every} lane that fired
anywhere along $w$.
\end{corollary}

\begin{proof}
A join in a join-semilattice equals the bottom iff every joinand is the bottom
(Proposition~\ref{prop:monoid}); apply to $\Phi_o(w)=\bigvee_i\varphi_o(s_i)$. The support of
the aggregate is $\supp\Phi_o(w)=\bigcup_i\supp\varphi_o(s_i)$, the union of the per-stage fired
lanes. \end{proof}

Corollary~\ref{cor:allstages} is the correctness statement for \code{ReceiptMeta.denial}: when
a manufacture is receipted, the honest denial word written into the \code{OcelCausalFrame} is
the aggregate $\Phi_o(w)$, so the receipt records exactly what was waved through --- not a
fiction of clean admission when a non-fatal lane fired. Composition (D1) is thereby a theorem,
not a convention.

\section{The fleet sweep is word-parallel}

\begin{proposition}[Byte-packed fleet admission]\label{prop:fleet}
Let a fleet of $N$ manufactures produce aggregate words $\Phi^{(1)},\dots,\Phi^{(N)}$, and pack
each into a byte $b^{(r)}=\firedmap(\Phi^{(r)})\in\{0,1\}^{8}$. Then:
\begin{enumerate}[label=(\roman*),leftmargin=*]
\item $\firedmap(\Phi_o(w))=\bigvee_{i}\firedmap(\varphi_o(s_i))$: the fleet byte of a pipeline
equals the bitwise \textsc{or} of its stage bytes (fired-mask commutes with aggregation).
\item A fleet admission sweep --- ``are all agents clean?'' --- is one linear pass computing
$\bigvee_r b^{(r)}$, eight agents per $64$-bit word, and returns $\bm 0$ iff every agent is
admitted.
\end{enumerate}
\end{proposition}

\begin{proof}
(i) $\firedmap$ is a homomorphism (Theorem~\ref{thm:firedhom}) and $\Phi_o$ is a homomorphism
(Theorem~\ref{thm:freehom}); their composite carries \code{compose} to $\lor$, so
$\firedmap(\bigvee_i\varphi_o(s_i))=\bigvee_i\firedmap(\varphi_o(s_i))$. (ii) The pass folds
$\lor$ over $N$ bytes; by Corollary~\ref{cor:allstages} applied fleet-wide, the fold is $\bm 0$
iff each $b^{(r)}=\bm 0$ iff each aggregate is $\Adml$ iff every agent is admitted. Packing eight
bytes per machine word makes the fold word-parallel. \end{proof}

Proposition~\ref{prop:fleet} is the algebraic underpinning of the keystone's ``agent as eight
bits'' construction and its memory-bandwidth-bounded planetary sweep: the reason a fleet of
$10^{12}$ agents can be admission-swept by masked \textsc{or}s is that both aggregation
(over stages) and projection (to bits) are semilattice homomorphisms, so the cheap bit-level
operation computes the same verdict as the expensive per-stage evaluation.

\chapter{Proof-Carrying Admission via SLD-Resolution}
\label{ch:proof}

\section{The logical quarantine}

Some admissions are not Boolean batteries but \emph{logical queries}: ``does this atom follow
from the admitted facts and rules?'' \code{praxis} routes these to the \code{prolog8} kernel
(\code{run\_kernel\_query} in \code{ops.rs}). Admission here is derivation, and --- unlike a
bare Boolean gate --- it can carry a \emph{proof}. First the language must be quarantined.

\begin{definition}[Logic admission; the Rice quarantine for Horn logic]\label{def:logicadm}
\code{prolog8} admits a query, fact block, or rule only if it lies in a bounded, decidable
Horn fragment: arity $\le 8$, rule body $\le 8$ atoms, $\le 8$ variables, proof fan-in $\le 8$;
no cut, no dynamic mutation (\code{assert}/\code{retract}), stratified negation, bounded or
declared recursion, no runtime text parsing, no side effects, interned (non-string) terms.
Violations are enumerated by \code{RejectionCode} (a closed enum of $\sim\!30$ structured
codes: \code{ArityCapExceeded}, \code{RuleBodyCapExceeded}, \code{CutNotAdmitted},
\code{UnstratifiedNegation}, \code{DynamicMutationNotAdmitted}, \dots). \code{admit\_atom} and
\code{admit\_rule} run before any query touches the kernel.
\end{definition}

Definition~\ref{def:logicadm} is Rice's theorem \cite{rice1953} answered at the logic layer:
full first-order or unrestricted Prolog derivability is undecidable, so the kernel retracts
onto a fragment whose byte caps make the search space finite, and everything outside the
fragment is refused with a code. This is the same move as $\adm:\Obs\to\Adm\cup\{\Rfsl\}$
(Paper~00), specialized to logic and carrying, additionally, a proof on the inside.

\section{The proof-carrying trichotomy}

\begin{definition}[Query result]\label{def:queryresult}
\code{Kernel::query} returns \code{QueryResult}, one of: \code{Answered(Vec<Decision>)},
\code{Denied(Box<Decision>)}, or \code{Invalid(RejectionCode)}. A \code{Decision} carries a
\emph{proof} (a \code{Vec<ProofNode>} DAG, positive or negative) and a \emph{receipt} sufficient
for deterministic replay.
\end{definition}

\begin{theorem}[Proof-carrying admission trichotomy]\label{thm:trichotomy}
For a query $q$ evaluated against an admitted kernel (admitted catalog, fact blocks, and
rules), exactly one of the following holds:
\begin{enumerate}[label=(\alph*),leftmargin=*]
\item \textbf{Answered.} $q$ passes \code{admit\_atom} and unifies with a stored fact or is
derivable by a depth-one rule application; then $\code{query}(q)=\code{Answered}(D_1,\dots,D_p)$
with $p\ge 1$, each $D_j$ carrying a \emph{positive} proof DAG whose every node has fan-in
$\le 8$, plus a receipt.
\item \textbf{Denied.} $q$ passes \code{admit\_atom} but no fact and no admitted rule derive it;
then $\code{query}(q)=\code{Denied}(D)$, where $D$ carries a \emph{negative} proof --- a finite,
closed witness that the depth-bounded SLD search was exhausted without success --- plus a
receipt.
\item \textbf{Invalid.} $q$ fails \code{admit\_atom}; then
$\code{query}(q)=\code{Invalid}(c)$ with $c\in\code{RejectionCode}$, and \emph{no proof is
produced}.
\end{enumerate}
Thus a \code{Denied} query yields a proof tree; an \code{Invalid} query yields a rejection code.
\end{theorem}

\begin{proof}
Trace \code{Kernel::query}. Its first action is the admission gate: if
\code{admit\_atom}$(q)$ returns \code{Err}$(c)$ the kernel returns \code{Invalid}$(c)$
immediately, before any search, establishing (c) and its no-proof clause. Otherwise the query
is in the admitted fragment and the kernel runs three phases. \emph{Phase 1} scans fact blocks
for rows unifying with $q$ under its bound positions; a nonempty match set is assembled into
\code{Answered} decisions, each with a positive proof rooted at the matched fact hash.
\emph{Phase 2}, on empty Phase 1, applies each rule whose head unifies with $q$ by depth-one
backward chaining with backtracking over body atoms (a single SLD-resolution step against the
fact base); any successful derivation is assembled into \code{Answered}, giving (a). The
fan-in bound holds because \code{admit\_rule} rejects any rule or proof node exceeding fan-in
$8$ (\code{ProofFanInExceeded}), so every emitted proof node has $\le 8$ children.
\emph{Phase 3} is reached iff both searches are empty; the kernel then assembles a
\code{Denied} decision with a \emph{negative} proof recording the exhausted search, giving (b).
The three phases are mutually exclusive by construction (each returns), so exactly one branch
fires. Soundness and completeness of the depth-bounded resolution --- that Phase 3 is reached
\emph{only} when $q$ is underivable within the fragment --- is SLD soundness and completeness
for definite/stratified Horn programs \cite{kowalski1974,lloyd1987}, restricted to the finite
search space that the byte caps of Definition~\ref{def:logicadm} guarantee. \end{proof}

\begin{proposition}[Boundedness gives a bounded certificate]\label{prop:boundedcert}
Every proof node emitted by the kernel has fan-in $\le 8$ and every rule body has $\le 8$
atoms, so a proof (positive or negative) is a bounded-degree DAG; its rolling
\code{proof\_root} is a single fixed-width hash field of the receipt. A bounded-context
verifier checks a query outcome by recomputing \code{proof\_root} and comparing --- cost
$O(\CL)$, independent of the size of the fact base --- and the deterministic replay path
recomputes exactly this root.
\end{proposition}

\begin{proof}
Fan-in and body caps are enforced by \code{admit\_rule} (\code{ProofFanInExceeded},
\code{RuleBodyCapExceeded}), so proof nodes have bounded out-degree; the receipt commits the
proof by its root hash (a projection of the DAG onto one field), and replay recomputes the
root from the stored decision, agreeing iff untampered --- the Faithful Projection mechanism of
the keystone applied to the proof DAG. The verifier reads the fixed-width root, not the DAG, so
its cost is $O(\CL)$. \end{proof}

Proposition~\ref{prop:boundedcert} is the projection principle at the logic layer: an
arbitrarily large search over facts and rules projects onto one root hash a bounded verifier
can check, so proof-carrying admission scales the same way byte-level admission does.

\section{Logic refusals rejoin the monoid}

The logical layer does not live apart from Chapters~\ref{ch:sl}--\ref{ch:pipe}; its refusals
are ordinary denial words.

\begin{proposition}[Kernel outcomes embed into the denial monoid]\label{prop:embed}
The refusing branches of Theorem~\ref{thm:trichotomy} embed into $(\Deny,\code{compose},\Adml)$
via the taxonomy: \code{run\_kernel\_query} maps a \code{Denied} outcome to the scenario
\code{KernelDenied} (category \code{Authorization}) and an \code{Invalid} outcome to
\code{KernelInvalid} (category \code{Identity}); by \code{denial\_lane} both compose into the
lane \code{CONFORMANCE\_GATE\_FAILED}. Hence a logical refusal is an element of $\Deny$ subject
to monotonicity (Theorem~\ref{thm:mono}), classification (Theorem~\ref{thm:total}), pipeline
composition (Theorem~\ref{thm:freehom}), and fired-mask projection (Theorem~\ref{thm:firedhom}),
exactly as any obligation refusal.
\end{proposition}

\begin{proof}
\code{run\_kernel\_query} constructs \code{RefusalScenario::KernelDenied} on the \code{Denied}
branch and \code{RefusalScenario::KernelInvalid} on the \code{Invalid} branch
(Definition~\ref{def:tax}); their categories are as stated by Theorem~\ref{thm:total}, and
\code{denial\_lane} sends both to \code{CONFORMANCE\_GATE\_FAILED}
(Proposition~\ref{prop:section}, closest-fit clause). Being elements of $\Deny$, they are acted
on by all the structure of Chapters~\ref{ch:sl}--\ref{ch:pipe}. \end{proof}

\begin{remark}[The epistemic distinction, made precise]
The trichotomy separates three epistemic states. \code{Invalid} means \emph{outside the
language}: a reason (a \code{RejectionCode}) but no proof --- the query was never well-formed
enough to have a truth value in the fragment. \code{Denied} means \emph{in the language and
false}: a negative proof witnesses underivability. \code{Answered} means \emph{in the language
and true}: a positive proof witnesses derivability. Only membership-in-the-language is a
type-like gate (\code{admit\_atom}); truth or falsity within the language is decided by bounded
search and \emph{always} ships a witness. This is the sharpest form of ``refusal is data'': at
the logic layer even a denial carries its proof, so a bounded verifier can trust the ``no''
without re-running the search.
\end{remark}

\chapter{Conclusion}
\label{ch:concl}

We set out to give the refusal object the algebra its role in a trillion-agent fleet demands,
and we have discharged each of the four demands of Chapter~\ref{ch:intro} as a theorem anchored
to \code{praxis}.

\emph{Composition (D1)} is Theorem~\ref{thm:freehom}: pipeline denial is the unique monoid
homomorphism from the free monoid on stages into the bounded commutative idempotent monoid
$(\Deny,\code{compose},\Adml)$ (Propositions~\ref{prop:monoid}, \ref{prop:code_is_sl}), so
aggregation is associative, order-independent, and multiplicity-independent, and a pipeline
admits iff every stage admits (Corollary~\ref{cor:allstages}). \emph{Classification (D2)} is the
category map $\catop$, and its \emph{totality (D3)} is Theorem~\ref{thm:total}: a
compiler-certified total classification onto exactly seven inhabited buckets plus a provably
empty \code{Reserved} slack coordinate, assembled with the total obligation and lane maps into
Corollary~\ref{cor:assembled} --- no refusal is unclassified. We located the exact boundary of
the type system's guarantee: the lane decoder is partial precisely because its domain is an
open $2^{64}$ newtype, so its exhaustiveness is a test obligation, not a type obligation
(Proposition~\ref{prop:partial}). And \emph{reason-carrying (D4)}, at its logical limit, is the
proof-carrying trichotomy of Theorem~\ref{thm:trichotomy}: an admitted-but-false query ships a
negative proof tree and an ill-formed query ships a rejection code, both bounded
(Proposition~\ref{prop:boundedcert}) and both rejoining the denial monoid
(Proposition~\ref{prop:embed}).

Monotonicity (Theorem~\ref{thm:mono}) ties the whole together: enlarging obligations only
tightens the gate, so the algebra is safe to grow. The fired-mask homomorphism
(Theorem~\ref{thm:firedhom}, Proposition~\ref{prop:fleet}) is what lets the keystone's
planetary admission sweep compute the true verdict with a bitwise \textsc{or}. What began, in
the keystone, as one Construction and one Proposition is now a complete small theory: the space
of ``why not'' is a join-semilattice, its classification is a mechanized surjection, its
aggregation is a free-monoid homomorphism, and its logical instance carries proof.

\begin{center}\itshape A refusal is not the absence of an answer; it is an answer with an algebra.\end{center}

\appendix
\chapter{Notation and Code Anchors}

\section*{Symbols}
\begin{center}
\begin{tabular}{ll}
\toprule
Symbol & Meaning\\
\midrule
$\Deny,\ \lor,\ \Adml$ & denial-word monoid; join (\code{compose}); bottom (\code{ADMITTED})\\
$\Deny_n,\ \supp d$ & abstract $n$-lane word space; fired-lane support\\
$\firedmap$ & fired-mask projection $\Deny\to\{0,1\}^{8}$ (\code{to\_fired\_mask})\\
$\delta_g,\ d_G(o)$ & obligation lane map; total denial (join over $G$)\\
$\Adm_G$ & admitted set under obligation set $G$\\
$\Catset,\ \Scnset$ & category set (8); scenario set (13)\\
$\catop,\ \scnop,\ \laneop,\ \scnop_\ell$ & category map; obligation$\to$scenario; scenario$\to$lane; lane decoder\\
$\Phi_o,\ \mathsf{Stage}^{*}$ & pipeline aggregate denial; free monoid on stages\\
$\Rfsl,\ \CL$ & refusal value; bounded-verifier context capacity\\
\bottomrule
\end{tabular}
\end{center}

\section*{Mathematics to code}
\begin{center}
\begin{tabular}{lll}
\toprule
Object & Code artifact & Location\\
\midrule
$(\Deny,\lor,\Adml)$ & \code{DenialPolarity}, \code{compose} & \code{bcinr-powl-receipt/src/denial.rs}\\
$\firedmap$ & \code{DenialPolarity::to\_fired\_mask} & \code{bcinr-powl-receipt/src/denial.rs}\\
$d_G$, \code{compose\_denials} & \code{compose\_denials} & \code{crates/praxis-core/src/refusal.rs}\\
$\catop$ (Thm~\ref{thm:total}) & \code{RefusalScenario::category} & \code{crates/praxis-core/src/refusal.rs}\\
$\scnop$ (Prop~\ref{prop:obtotal}) & \code{From<\&Obligation>} & \code{crates/praxis-core/src/refusal.rs}\\
$\scnop_\ell$ (Prop~\ref{prop:partial}) & \code{scenario\_for\_denial\_lane} & \code{crates/praxis-core/src/refusal.rs}\\
$\laneop$ (Prop~\ref{prop:section}) & \code{denial\_lane} & \code{crates/praxis-core/src/refusal.rs}\\
\code{Obligation} (3 kinds) & \code{Obligation} enum & \code{crates/praxis-core/src/law.rs}\\
\code{Andon::Halted} & \code{Andon} enum & \code{crates/praxis-core/src/law.rs}\\
Trichotomy (Thm~\ref{thm:trichotomy}) & \code{Kernel::query}, \code{QueryResult} & \code{prolog8/src/kernel.rs}\\
Logic quarantine & \code{admit\_atom}, \code{RejectionCode} & \code{prolog8/src/admission.rs}\\
Kernel refusals (Prop~\ref{prop:embed}) & \code{run\_kernel\_query} & \code{src/ops.rs}\\
\bottomrule
\end{tabular}
\end{center}

\begin{thebibliography}{9}

\bibitem{chatman2025}
S.~Chatman,
\emph{The Chatman Equation and the Industrial Revolution of Knowledge:
$A=\mu(O)$, Knowledge Hooks, and Production-Verified Enterprise Execution},
v1.2.0, November 2025.

\bibitem{rice1953}
H.~G.~Rice,
\emph{Classes of recursively enumerable sets and their decision problems},
Transactions of the American Mathematical Society \textbf{74} (1953), 358--366.

\bibitem{birkhoff1967}
G.~Birkhoff,
\emph{Lattice Theory},
3rd ed., American Mathematical Society Colloquium Publications, vol.~25, 1967.

\bibitem{kowalski1974}
R.~A.~Kowalski,
\emph{Predicate logic as programming language},
Information Processing \textbf{74} (Proc.\ IFIP Congress), North-Holland, 1974, 569--574.

\bibitem{lloyd1987}
J.~W.~Lloyd,
\emph{Foundations of Logic Programming},
2nd ed., Springer-Verlag, 1987.

\bibitem{vdaalst2003patterns}
W.~M.~P.~van der Aalst, A.~H.~M.~ter Hofstede, B.~Kiepuszewski, and A.~P.~Barros,
\emph{Workflow patterns},
Distributed and Parallel Databases \textbf{14}(1) (2003), 5--51.

\bibitem{blake3}
J.~O'Connor, J.-P.~Aumasson, S.~Neves, and Z.~Wilcox-O'Hearn,
\emph{BLAKE3: One function, fast everywhere},
2020.

\end{thebibliography}

\end{document}
